\documentclass[../../main.tex]{subfiles}% 注意这里的文件路径不能用 ./main.tex ,否则用latexmk编译子文件会报错
\graphicspath{{\subfix{./image/}}{\subfix{../../image/}}} % 指定图片引用路径,后续可以直接使用图片文件名
% 如果图片存放在上级目录,则使用 ../../image/ 作为备用路径

% 例如:
% \begin{figure}[H]
% \centering
% \includegraphics[width=0.3\textwidth]{图.png}
% \caption{}
% \label{figure:图}
% \end{figure}
% 注意:上述\label{}一定要放在\caption{}之后,否则引用图片序号会只会显示??.

\begin{document}

\section{Young 测度再探}

设 $X$ 为局部紧 Hausdorff 空间，$\{u_k\}\subseteq L^1(X,\mathscr{B}(X),\mu)$ 弱收敛。已观察到（见\cref{example:ch02-弱收敛振荡}(i)），对非线性 $h$（使 $\{h\circ u_k\}\subseteq L^1(X,\mathscr{B}(X),\mu)$），$\{h\circ u_k\}$ 未必如此。为获得 $\{h\circ u_k\}$ 性态的信息，\cref{remark:ch05-函数转测度}建议我们处理与 $\{h\circ u_k\}$ 相伴的序列 $\{\nu_k\}\subseteq\mathscr{R}_f(X)$，如\eqref{align:ch05-5.5}。在此框架中，与 $\{u_k\}$ 相伴的 Young 测度序列 $\{\nu_{u_k}\}$ 起重要作用。

设 $\mu\in\mathscr{R}^+_f(X)$，$\nu\in\mathcal{Y}^+(X,\mathscr{B}(X),\mu;Y)\equiv\mathcal{Y}^+(X;Y)$ 为 Young 测度。由定义 $\nu_{p_X}=\mu$（见\cref{definition:ch02-Young测度}），且由于 $\mathcal{Y}^+(X;Y)\subseteq\mathscr{R}^+_f(X\times Y)$，由\cref{theorem:ch05-离散分解}得如下结果。

\begin{proposition}[Young 测度的离散分解]\label{proposition:ch05-离散分解}
设 $X$ 与 $Y$ 为具可数基的局部紧 Hausdorff 空间。设 $\mu\in\mathscr{R}^+_f(X)$，$\nu\in\mathcal{Y}^+(X;Y)$。则存在 $Y$ 上的参数化测度 $\nu\equiv\{\nu_x\}_{x\in X}$ 具有下列性质：
\begin{enumerate}[(i)]
\item\label{item:ch05-youdis-i} $\nu_x\in P(Y)\subseteq\mathscr{R}^+_f(Y)$ 对 $\mu$-几乎所有 $x\in X$；
\item\label{item:ch05-youdis-ii} 对每个 $h\in L^1(X\times Y,\mathscr{B}(X\times Y),\nu)$，
\begin{align}
h(x,\cdot)\in L^1(Y,\mathscr{B}(Y),\nu_x)
\quad\text{对 }\mu\text{-几乎所有 }x\in X,\label{align:ch05-5.48}\\
\text{映射 }x\mapsto\int_Y h(x,y)\,\mathrm{d}\nu_x(y)
\text{ 属于 }L^1(X,\mathscr{B}(X),\mu),\label{align:ch05-5.49}
\end{align}
且
\begin{align}
\int_{X\times Y}h(x,y)\,\mathrm{d}\nu(x,y)
=\int_X\mathrm{d}\mu(x)\int_Y h(x,y)\,\mathrm{d}\nu_x(y).\label{align:ch05-5.50}
\end{align}
\end{enumerate}
若 $\{\nu'_x\}_{x\in X}$ 为 $Y$ 上的参数化测度使
\begin{align}
\text{映射 }x\mapsto\int_Y h(x,y)\,\mathrm{d}\nu'_x(y)
\text{ 属于 }L^1(X,\mathscr{B}(X),\mu)\label{align:ch05-5.51}
\end{align}
且
\begin{align}
\int_{X\times Y}h(x,y)\,\mathrm{d}\nu(x,y)
=\int_X\mathrm{d}\mu(x)\int_Y h(x,y)\,\mathrm{d}\nu'_x(y)\label{align:ch05-5.52}
\end{align}
对每个有界 $\mathscr{B}(X\times Y)$-可测、紧支撑的 $h:X\times Y\to\mathbb{R}$ 成立，则 $\nu'_x=\nu_x$ 对 $\mu$-几乎所有 $x\in X$。
\end{proposition}

刻画与函数相伴的 Young 测度的离散分解是有趣的（见\cref{definition:ch02-Young测度}）。回忆与 $u$ 相伴的 Young 测度 $\nu_u$ 定义为
\begin{align*}
\nu_u(E\times F)\triangleq \mu(E\cap u^{-1}(F))
\quad\text{对所有 }E\in\mathscr{B}(X),\ F\in\mathscr{B}(Y).
\end{align*}

\begin{proposition}\label{proposition:ch05-Dirac测度}
设 $X$ 与 $Y$ 为具可数基的局部紧 Hausdorff 空间。设 $\mu\in\mathscr{R}^+_f(X)$，$u:X\to Y$ 为 $(\mathscr{B}(X),\mathscr{B}(Y))$-可测，$\nu_u\in\mathcal{Y}^+(X;Y)$ 为与 $u$ 相伴的 Young 测度。则
\begin{align}
(\nu_u)_x=\delta_{u(x)}
\quad\text{对 }\mu\text{-几乎所有 }x\in X.\label{align:ch05-5.53}
\end{align}
\end{proposition}

\begin{proof}
对任意 $F\in\mathscr{B}(Y)$，由\eqref{align:ch05-5.36}得
\begin{align*}
\nu_u(X\times F)=\int_{X\times Y}\chi_F(y)\,\mathrm{d}\nu_u(x,y)
=\int_X(\nu_u)_x(F)\,\mathrm{d}\mu(x).
\end{align*}
另一方面，由定义
\begin{align*}
\nu_u(X\times F)=\mu(u^{-1}(F))
=\int_X\chi_{u^{-1}(F)}(x)\,\mathrm{d}\mu(x).
\end{align*}
由\eqref{align:ch02-2.29}，对每个 $F\in\mathscr{B}(Y)$，存在 $\mu$-零测集 $N\subseteq X$ 使对所有 $x\in N^c$，
\begin{align*}
(\nu_u)_x(F)=\chi_{u^{-1}(F)}(x)=
\begin{cases}
1 & \text{若 }u(x)\in F,\\
0 & \text{否则}.
\end{cases}
\end{align*}
由于 $X$ 具有可数基，可测空间 $(X,\mathscr{B}(X))$ 可分（见\cref{definition:ch01-可分空间}），故由标准论证，$\mu$-零测集 $N$ 的选取可独立于 $F\in\mathscr{B}(Y)$。结论得证。
\end{proof}

由等式\eqref{align:ch05-5.36}与\eqref{align:ch05-5.53}立即得如下结果。

\begin{corollary}\label{corollary:ch05-Young积分}
设 $X$ 与 $Y$ 为具可数基的局部紧 Hausdorff 空间。设 $\mu\in\mathscr{R}^+_f(X)$，$u:X\to Y$ 为 $(\mathscr{B}(X),\mathscr{B}(Y))$-可测，$\nu_u\in\mathcal{Y}^+(X,\mathscr{B}(X),\mu;Y)$ 为与 $u$ 相伴的 Young 测度。则对每个 $h\in L^1(X\times Y,\mathscr{B}(X\times Y),\nu_u)$，
\begin{align}
\int_{X\times Y}h(x,y)\,\mathrm{d}\nu_u(x,y)
=\int_X h(x,u(x))\,\mathrm{d}\mu(x).\label{align:ch05-5.54}
\end{align}
\end{corollary}

鉴于\cref{proposition:ch05-离散分解}，考虑从 $X$ 到 $\mathscr{R}_f(Y)$ 的弱$^*$可测函数集是有趣的，这些函数对 $\mu$-几乎所有 $x\in X$ 取 $Y$ 上概率测度集 $P(Y)\subseteq\mathscr{R}^+_f(Y)$ 中的值。我们记 $P(X;Y)\equiv P(X,\mathscr{B}(X),\mu;Y)$ 为该集，称为 $Y$ 上的\textbf{参数化概率}集。易见 $P(X;Y)$ 是空间 $L^\infty_{w^*}(X;\mathscr{R}_f(Y))$ 的单位球
\begin{align}
U\triangleq \left\{\nu\in L^\infty_{w^*}(X;\mathscr{R}_f(Y))\ \middle|\
\operatorname*{ess\,sup}_{x\in X}\|\nu_x\|_{\mathscr{R}_f(Y)}\leqslant1\right\}\label{align:ch05-5.55}
\end{align}
的闭凸子集。

\begin{corollary}\label{corollary:ch05-Young一一对应}
设 $X$ 与 $Y$ 为具可数基的局部紧 Hausdorff 空间，$\mu\in\mathscr{R}^+_f(X)$。则 $\mathcal{Y}^+(X;Y)$ 与 $P(X;Y)$ 之间存在一一对应。
\end{corollary}

\begin{proof}
由\cref{proposition:ch05-离散分解}，每个 $\nu\in\mathcal{Y}^+(X;Y)$ 对应 $Y$ 上的一个参数化概率。反之，设 $\{\nu_x\}_{x\in X}\in P(X;Y)$，置
\begin{align*}
\nu:\mathscr{B}(X\times Y)\to[0,\infty),\quad
\nu(G)\triangleq \int_X\nu_x(G_x)\,\mathrm{d}\mu(x)
\quad(G\in\mathscr{B}(X\times Y))
\end{align*}
（注意由\cref{remark:ch05-参数测度等价}该定义良定）。由\cref{proposition:ch01-Radon投影}，$\nu\in\mathscr{R}^+_f(X\times Y)$，取 $G=E\times Y$（$E\in\mathscr{B}(X)$）得
\begin{align*}
\nu(E\times Y)=\int_E\nu_x(Y)\,\mathrm{d}\mu(x)=\mu(E).
\end{align*}
故 $\nu\in\mathcal{Y}^+(X;Y)$（见\cref{definition:ch02-Young测度}），结论得证。
\end{proof}

由\cref{corollary:ch05-Young一一对应}与前述注记，自然赋予 $\mathcal{Y}^+(X;Y)$ 以 $L^\infty_{w^*}(X;\mathscr{R}_f(Y))$ 上弱$^*$收敛的拓扑（见\cref{subsection:ch05-Young弱星}）。$\mathcal{Y}^+(X;Y)$ 上另一种与窄收敛相关的拓扑将在\cref{subsection:ch05-Young窄}中引入。为此，我们需要如下定义。

\begin{definition}\label{definition:ch05-Caratheodory}
设 $(X,\mathcal{A},\mu)$ 为有限测度空间，$Y$ 为拓扑空间。
\begin{enumerate}[(i)]
\item\label{item:ch05-car-i} \textbf{Carathéodory 函数}是指任意 $h:X\times Y\to\mathbb{R}$ 满足
\begin{align*}
\begin{cases}
\text{(a) } h(x,\cdot)\in C(Y) & \text{对 }\mu\text{-几乎所有 }x\in X;\\
\text{(b) } h(\cdot,y)\text{ 为 }\mathcal{A}\text{-可测} & \text{对所有 }y\in Y.
\end{cases}
\end{align*}
\item\label{item:ch05-car-ii} \textbf{有界 Carathéodory 被积函数}是指任意 Carathéodory 函数满足
\begin{align*}
\begin{cases}
\text{(a) } h(x,\cdot)\in C_b(Y) & \text{对 }\mu\text{-几乎所有 }x\in X;\\
\text{(b) 映射 }x\mapsto\|h(x,\cdot)\|_\infty\text{ 为 }\mathcal{A}\text{-可测},
\text{ 且 }\int_X\|h(x,\cdot)\|_\infty\,\mathrm{d}\mu(x)<\infty.
\end{cases}
\end{align*}
\end{enumerate}
$C_0$-Carathéodory 被积函数由在 (a) 中将 $C_b(Y)$ 替换为 $C_0(Y)$ 定义。

两个 Carathéodory 函数 $h$ 与 $h'$ 属于同一等价类，若 $h(x,\cdot)=h'(x,\cdot)$ 对 $\mu$-几乎所有 $x\in X$ 成立。
\end{definition}

Carathéodory 函数（等价类）所成的向量空间记作 $\mathcal{C}(X\times Y)$，有界 Carathéodory 被积函数与 $C_0$-Carathéodory 被积函数所成的空间分别记作 $\mathcal{C}_b(X\times Y)$ 与 $\mathcal{C}_0(X\times Y)$。

立即可见，映射 $\mathcal{C}_b(X\times Y)\to[0,\infty)$，
\begin{align}
h\mapsto\|h\|_{\mathcal{C}_b(X\times Y)}
\triangleq \int_X\|h(x,\cdot)\|_\infty\,\mathrm{d}\mu(x),\label{align:ch05-5.56}
\end{align}
是一个范数。于是我们有如下结果。

\begin{proposition}\label{proposition:ch05-Caratheodory同构}
设 $(X,\mathcal{A},\mu)$ 为有限测度空间，$Y$ 为具可数基的局部紧 Hausdorff 空间。则映射
\begin{align}
T:\mathcal{C}_0(X\times Y)\to L^1(X;C_0(Y)),\qquad
(Th)(x)\triangleq h(x,\cdot)
\quad\text{对 }\mu\text{-几乎所有 }x\in X,\label{align:ch05-5.57}
\end{align}
是 $(\mathcal{C}_0(X\times Y),\|\cdot\|_{\mathcal{C}_0(X\times Y)})$ 到 $L^1(X;C_0(Y))$ 的等距同构。特别地，$\mathcal{C}_0(X\times Y)$ 赋以范数 $\|\cdot\|_{\mathcal{C}_b(X\times Y)}$ 为 Banach 空间。
\end{proposition}

\begin{proof}
我们先证算子 $T$ 良定。对任意 $h\in\mathcal{C}_0(X\times Y)$ 与 $g\in C_0(Y)$，函数 $(x,y)\mapsto h(x,y)-g(y)$ 属于 $\mathcal{C}_0(X\times Y)$，故对每个 $\alpha\in\mathbb{R}$，集合 $\{x\in X\mid\|h(x,\cdot)-g\|_\infty<\alpha\}$ 属于 $\mathcal{A}$（见\cref{definition:ch05-Caratheodory}(ii)）。另一方面，$C_0(Y)$ 可分（见附录 A,\cref{subsection:appA-可分}），故 Borel $\sigma$-代数 $\mathscr{B}\equiv\mathscr{B}(C_0(Y))$ 由可数开球基 $B(g_k,\alpha)\triangleq \{g\in C_0(Y)\mid\|g-g_k\|_\infty<\alpha\}$（$k\in\mathbb{N}$）生成。由此映射 $x\mapsto h(x,\cdot)=(Th)(x)$ 为 $(\mathcal{A},\mathscr{B})$-可测，因而由\cref{proposition:ch04-可测性}为 $\mu$-可测，因为 $C_0(Y)$ 可分。此外，
\begin{align}
\|h\|_{\mathcal{C}_b(X\times Y)}
=\int_X\|h(x,\cdot)\|_\infty\,\mathrm{d}\mu(x)
=\int_X\|(Th)(x)\|_\infty\,\mathrm{d}\mu(x)
=\|Th\|_{L^1(X;C_0(Y))}<\infty.\label{align:ch05-5.58}
\end{align}
故 $Th\in L^1(X;C_0(Y))$，断言得证。

由于 $T$ 线性且等距（见\eqref{align:ch05-5.58}），它单射。为证满射，固定任意 $\hat{h}\in L^1(X;C_0(Y))$。则对 $\mu$-几乎所有 $x\in X$，$\hat{h}(x)\in C_0(Y)$，且映射 $x\mapsto\|\hat{h}(x)\|_\infty$ 属于 $L^1(X)$。此外，对任意 $y\in Y$，映射 $x\mapsto\hat{h}(x)(y)$ 为 $\mathcal{A}$-可测，因为映射 $x\mapsto\|\hat{h}(x)\|_\infty$ 为 $\mathcal{A}$-可测，且映射 $C_0(Y)\to\mathbb{R}$，$g\mapsto g(y)$（$g\in C_0(Y)$，$y\in Y$）连续。因此函数 $(x,y)\mapsto\hat{h}(x)(y)$ 属于 $\mathcal{C}_0(X\times Y)$，故 $T$ 满射。结论得证。
\end{proof}

由\cref{proposition:ch05-Caratheodory同构}、\cref{theorem:ch04-可分}与\cref{proposition:ch04-对偶}得如下结果。

\begin{corollary}\label{corollary:ch05-Caratheodory对偶}
设 $(X,\mathcal{A},\mu)$ 为有限测度空间，$Y$ 为具可数基的局部紧 Hausdorff 空间。则：
\begin{enumerate}[(i)]
\item\label{item:ch05-cd-i} 空间 $\mathcal{C}_0(X\times Y)$ 可分；
\item\label{item:ch05-cd-ii} 我们有
\begin{align*}
(\mathcal{C}_0(X\times Y))^*=L^\infty_{w^*}(X;\mathscr{R}_f(Y)).
\end{align*}
\end{enumerate}
\end{corollary}

\begin{remark}\label{remark:ch05-紧空间同构}
设 $(X,\mathcal{A},\mu)$ 为有限测度空间，$Y$ 为紧空间。则 $C_b(Y)$ 可分，且\cref{proposition:ch05-Caratheodory同构}的证明表明映射
\begin{align}
T:\mathcal{C}_b(X\times Y)\to L^1(X;C_b(Y)),\qquad
(Th)(x)\triangleq h(x,\cdot)
\quad\text{对 }\mu\text{-几乎所有 }x\in X,\label{align:ch05-5.59}
\end{align}
是 $(\mathcal{C}_b(X\times Y),\|\cdot\|_{\mathcal{C}_b(X\times Y)})$ 到 $L^1(X;C_b(Y))$ 的等距同构。
\end{remark}

\subsection{弱$^*$收敛}\label{subsection:ch05-Young弱星}

如已观察到的，由\cref{corollary:ch05-Young一一对应}与前述注记，自然赋予 $\mathcal{Y}^+(X;Y)$ 以 $L^\infty_{w^*}(X;\mathscr{R}_f(Y))$ 上弱$^*$收敛的拓扑。于是由\cref{corollary:ch05-Caratheodory对偶}(ii)得如下定义（回忆由定义 $\mathcal{Y}^+(X;Y)\subseteq\mathscr{R}^+_f(X\times Y)$）。

\begin{definition}\label{definition:ch05-Young弱星}
设 $X$ 与 $Y$ 为具可数基的局部紧 Hausdorff 空间。序列 $\{\nu_k\}\subseteq\mathcal{Y}^+(X;Y)$ \textbf{弱$^*$收敛}于 $\nu\in L^\infty_{w^*}(X;\mathscr{R}_f(Y))$（记作 $\nu_k\stackrel{*}{\rightharpoonup}\nu$），若
\begin{align*}
\lim_{k\to\infty}\int_{X\times Y}h\,\mathrm{d}\nu_k
=\int_{X\times Y}h\,\mathrm{d}\nu
\quad\text{对所有 }h\in\mathcal{C}_0(X\times Y).
\end{align*}
$\mathcal{Y}^+(X;Y)$ 上的弱$^*$拓扑，记作 $\mathcal{T}_{w^*}$，是使所有映射 $\nu\mapsto\int_{X\times Y}h\,\mathrm{d}\nu$（$h\in\mathcal{C}_0(X\times Y)$）连续的最弱拓扑。

$\mathcal{T}_{w^*}$ 的一组基为子集族
\begin{align}
B_{h_1,\cdots ,h_n;\epsilon}(\nu_0)
\triangleq \left\{\nu\in\mathcal{Y}^+(X;Y)\ \middle|\
\left|\int_{X\times Y}h_j\,\mathrm{d}\nu-\int_{X\times Y}h_j\,\mathrm{d}\nu_0\right|<\epsilon
\ \forall j=1,\cdots ,n\right\},\label{align:ch05-5.60}
\end{align}
其中 $h_1,\cdots ,h_n\in\mathcal{C}_0(X\times Y)$，$\nu_0\in\mathcal{Y}^+(X;Y)$，$n\in\mathbb{N}$，$\epsilon>0$。
\end{definition}

\begin{remark}\label{remark:ch05-弱星损失质量}
由\cref{definition:ch05-Young弱星}与 $L^\infty_{w^*}(X;\mathscr{R}_f(Y))$ 中范数的下半连续性，弱$^*$极限 $\nu$ 属于单位球 $U\subseteq L^\infty_{w^*}(X;\mathscr{R}_f(Y))$（见\eqref{align:ch05-5.55}）。然而，由于可能的“质量损失”，$\nu$ 不必是 $Y$ 上的参数化概率（见\cref{remark:ch05-范数下半连续}）。还注意 $\nu\in L^\infty_{w^*}(X;\mathscr{R}^+_f(Y))$，因为 $\{\nu_k\}\subseteq\mathcal{Y}^+(X;Y)$。
\end{remark}

\begin{theorem}\label{theorem:ch05-Young弱星紧}
设 $X$ 与 $Y$ 为具可数基的局部紧 Hausdorff 空间。则对每个序列 $\{\nu_k\}\subseteq\mathcal{Y}^+(X;Y)$，存在子列 $\{\nu_{k_l}\}\subseteq\{\nu_k\}$ 与 $\nu\in U\subseteq L^\infty_{w^*}(X;\mathscr{R}_f(Y))$ 使 $\nu_{k_l}\stackrel{*}{\rightharpoonup}\nu$。
\end{theorem}

\begin{proof}
由\cref{theorem:ch04-可分}，空间 $L^1(X;C_0(Y))$ 可分，故由\cref{proposition:ch04-对偶}与 Banach 定理，单位球 $U\subseteq L^\infty_{w^*}(X;\mathscr{R}_f(Y))$ 弱$^*$紧。于是由\cref{corollary:ch05-Caratheodory对偶}结论得证。
\end{proof}

\begin{remark}\label{remark:ch05-Young弱星度量}
由于 $L^1(X;C_0(Y))=\mathcal{C}_0(X\times Y)$ 可分，相对弱$^*$拓扑 $\mathcal{T}_{w^*}\cap U$ 可度量化。事实上，由\cref{proposition:appA-弱星可度量化}，它与由范数
\begin{align}
\nu\mapsto|||\nu|||
\triangleq \sum_{k=1}^\infty\frac{1}{2^k}
\frac{\left|\int_{X\times Y}h_k\,\mathrm{d}\nu\right|}
{1+\int_X\|h_k(x,\cdot)\|_\infty\,\mathrm{d}\mu(x)}\label{align:ch05-5.61}
\end{align}
在 $L^\infty_{w^*}(X;\mathscr{R}_f(Y))$ 上诱导的度量拓扑一致，其中 $D\equiv\{h_k\}\subseteq\mathcal{C}_0(X\times Y)$ 为可数稠密子集。
\end{remark}

\subsection{窄收敛与胎紧性}\label{subsection:ch05-Young窄}

\begin{definition}\label{definition:ch05-Young窄}
设 $X$ 与 $Y$ 为具可数基的局部紧 Hausdorff 空间。序列 $\{\nu_k\}\subseteq\mathcal{Y}^+(X;Y)$ \textbf{窄收敛}于 $\nu\in\mathcal{Y}^+(X;Y)$（记作 $\nu_k\stackrel{n}{\rightharpoonup}\nu$），若
\begin{align*}
\lim_{k\to\infty}\int_{X\times Y}h\,\mathrm{d}\nu_k
=\int_{X\times Y}h\,\mathrm{d}\nu
\quad\text{对所有 }h\in\mathcal{C}_b(X\times Y).
\end{align*}
$\mathcal{Y}^+(X;Y)$ 上的窄拓扑，记作 $\mathcal{T}_n$，是使所有映射 $\nu\mapsto\int_{X\times Y}h\,\mathrm{d}\nu$（$h\in\mathcal{C}_b(X\times Y)$）连续的最弱拓扑。

$\mathcal{T}_n$ 的一组基为子集族\eqref{align:ch05-5.60}，只是其中 $h_1,\cdots ,h_n\in\mathcal{C}_b(X\times Y)$。
\end{definition}

如下半连续性结果将被使用。

\begin{proposition}\label{proposition:ch05-Caratheodory下半连续}
设 $X$ 与 $Y$ 为具可数基的局部紧 Hausdorff 空间，$h:X\times Y\to[0,\infty)$ 为 Carathéodory 函数。则映射 $\mathcal{Y}^+(X;Y)\to[0,\infty]$，$\nu\mapsto\int_{X\times Y}h\,\mathrm{d}\nu$，在窄拓扑中下半连续。
\end{proposition}

\begin{proof}
我们须证明：若在 $\mathcal{Y}^+(X;Y)$ 中 $\nu_k\stackrel{n}{\rightharpoonup}\nu$，则
\begin{align}
\int_{X\times Y}h(x,y)\,\mathrm{d}\nu(x,y)
\leqslant\liminf_{k\to\infty}\int_X\mathrm{d}\mu(x)\int_Y h(x,y)\,\mathrm{d}(\nu_k)_x(y).\label{align:ch05-5.62}
\end{align}
设 $d$ 为空间 $Y$ 上的相容度量（由\cref{proposition:appA-紧生成}，$Y$ 可度量化）。对每个 $n\in\mathbb{N}$，定义 $h_n:X\times Y\to[0,\infty)$ 为
\begin{align*}
h_n(x,y)\triangleq \min\left\{\inf_{\zeta\in Y}[h(x,\zeta)+n\,d(\zeta,y)],\ n\right\}
\qquad(x\in X,\ y\in Y).
\end{align*}
易见对每个 $n\in\mathbb{N}$，$h_n$ 为 $\mathscr{B}(X\times Y)$-可测，$h_n\in\mathcal{C}_b(X\times Y)$，$h_n\leqslant h_{n+1}$，且 $h_n\to h$ 在 $X\times Y$ 上逐点收敛。于是由单调收敛定理，
\begin{align*}
\int_{X\times Y}h\,\mathrm{d}\nu
=\lim_{n\to\infty}\int_{X\times Y}h_n\,\mathrm{d}\nu
=\sup_{n\in\mathbb{N}}\int_{X\times Y}h_n\,\mathrm{d}\nu.
\end{align*}
对每个 $n\in\mathbb{N}$，映射 $\nu\mapsto\int_{X\times Y}h_n\,\mathrm{d}\nu$ 由\cref{definition:ch05-Young窄}连续，因为 $h_n\in\mathcal{C}_b(X\times Y)$。故映射 $\nu\mapsto\int_{X\times Y}h\,\mathrm{d}\nu$ 作为连续函数序列的上确界为下半连续。结果得证。
\end{proof}

对 Young 测度，胎紧性定义如下。

\begin{definition}\label{definition:ch05-Young胎紧}
设 $X$ 与 $Y$ 为具可数基的局部紧 Hausdorff 空间。子集 $N\subseteq\mathcal{Y}^+(X;Y)$ 称为\textbf{胎紧的}，若对任意 $\epsilon>0$，存在紧子集 $K\subseteq Y$ 使对所有 $\nu\in N$ 有 $\nu(X\times K^c)<\epsilon$。
\end{definition}

\begin{remark}\label{remark:ch05-Young胎紧投影}
由\cref{definition:ch05-Young胎紧}，子集 $N\subseteq\mathcal{Y}^+(X;Y)$ 胎紧，当且仅当投影到 $Y$ 的集合 $\{\nu_{p_Y}\mid\nu\in N\}\subseteq\mathscr{R}^+_f(Y)$ 在\cref{definition:ch05-胎紧}意义下胎紧。
\end{remark}

如下结果的证明与\cref{proposition:ch05-胎紧拓扑一致}严格类似，故我们略去。

\begin{proposition}\label{proposition:ch05-Young胎紧一致}\label{proposition:ch05-矩估计}
设 $X$ 与 $Y$ 为具可数基的局部紧 Hausdorff 空间，$N\subseteq\mathcal{Y}^+(X;Y)$ 胎紧。则 $N$ 上的窄拓扑与 $N$ 上的弱$^*$拓扑一致，即 $\mathcal{T}_n\cap N=\mathcal{T}_{w^*}\cap N$。
\end{proposition}

作为胎紧性判据，由\cref{remark:ch05-Young胎紧投影}与\cref{proposition:ch05-胎紧判据}得如下结果。

\begin{proposition}\label{proposition:ch05-Young胎紧判据}
设 $X$ 与 $Y$ 为具可数基的局部紧 Hausdorff 空间。则下列命题等价：
\begin{enumerate}[(i)]
\item\label{item:ch05-ytight-i} $N\subseteq\mathcal{Y}^+(X;Y)$ 胎紧；
\item\label{item:ch05-ytight-ii} 存在下紧函数 $\varphi:Y\to[0,\infty)$ 使
\begin{align}
\sup_{\nu\in N}\int_Y\varphi\,\mathrm{d}\nu_{p_Y}<\infty.\label{align:ch05-5.63}
\end{align}
\end{enumerate}
\end{proposition}

对 Young 测度，我们有\cref{theorem:ch05-prokhorov}的如下对应物。

\begin{theorem}\label{theorem:ch05-弱紧Young}\label{proposition:ch05-弱紧YoungL1}
设 $X$ 与 $Y$ 为具可数基的局部紧 Hausdorff 空间。设 $N\subseteq\mathcal{Y}^+(X;Y)$。则下列命题等价：
\begin{enumerate}[(i)]
\item\label{item:ch05-youngcomp-i} $N$ 在窄拓扑中相对序列紧；
\item\label{item:ch05-youngcomp-ii} $N$ 胎紧。
\end{enumerate}
\end{theorem}

\begin{proof}
(i)$\Rightarrow$(ii)。由假设，对任意序列 $\{\nu_k\}\subseteq N$，存在子列 $\{\nu_{k_j}\}\subseteq\{\nu_k\}$ 与 $\nu\in\mathcal{Y}^+(X;Y)$ 使
\begin{align}
\lim_{j\to\infty}\int_{X\times Y}h\,\mathrm{d}\nu_{k_j}
=\int_{X\times Y}h\,\mathrm{d}\nu
\quad\text{对所有 }h\in\mathcal{C}_b(X\times Y).\label{align:ch05-5.64}
\end{align}
取 $h(x,y)=g(y)$，$g\in C_b(Y)$。由于
\begin{align*}
\int_{X\times Y}g(y)\,\mathrm{d}\nu(x,y)=\int_Y g(y)\,\mathrm{d}\nu_{p_Y}(y),\\
\int_{X\times Y}g(y)\,\mathrm{d}\nu_{k_j}(x,y)=\int_Y g(y)\,\mathrm{d}(\nu_{k_j})_{p_Y}(y)
\end{align*}
（见\cref{remark:ch02-Young投影}(i)），由\eqref{align:ch05-5.64}得
\begin{align*}
\lim_{j\to\infty}\int_Y g(y)\,\mathrm{d}(\nu_{k_j})_{p_Y}(y)
=\int_Y g(y)\,\mathrm{d}\nu_{p_Y}(y)
\quad\text{对所有 }g\in C_b(Y).
\end{align*}
由此集合 $\{\nu_{p_Y}\mid\nu\in N\}$ 在 $\mathscr{R}^+_f(Y)$ 的窄拓扑中相对序列紧，故由\cref{theorem:ch05-prokhorov}胎紧。于是由\cref{remark:ch05-Young胎紧投影}断言得证。

(ii)$\Rightarrow$(i)。由\cref{remark:ch02-Young投影}，$N$ 是 $\mathscr{R}^+_f(X\times Y)$ 的有界子集。我们证明它也是 $\mathscr{R}^+_f(X\times Y)$ 的胎紧子集，即对任意 $\epsilon>0$，存在紧子集 $\tilde{K}\subseteq X\times Y$ 使对所有 $\nu\in N$ 有 $\nu(\tilde{K}^c)<\epsilon$（见\cref{definition:ch05-胎紧}）。由假设，对任意 $\epsilon>0$，存在紧子集 $K\subseteq Y$ 使对所有 $\nu\in N$ 有 $\nu_{p_Y}(K^c)=\nu(X\times K^c)<\epsilon/2$。此外，由\cref{proposition:appA-紧生成}，$X$ 为 $\sigma$-紧，故存在紧子集 $K_0\subseteq X$ 使 $\mu(K_0^c)<\epsilon/2$。注意由 Young 测度定义，对每个 $\nu\in N$ 有 $\nu_{p_X}(K_0^c)=\nu(K_0^c\times Y)=\mu(K_0^c)$。置 $\tilde{K}\triangleq K_0\times K\subseteq X\times Y$。显然 $\tilde{K}$ 紧，且 $\tilde{K}^c\subseteq(X\times K^c)\cup(K_0^c\times Y)$。则对所有 $\nu\in N$，
\begin{align*}
\nu(\tilde{K}^c)\leqslant\nu(X\times K^c)+\nu(K_0^c\times Y)
=\nu_{p_Y}(K^c)+\mu(K_0^c)<\epsilon,
\end{align*}
断言得证。

由上述注记与\cref{proposition:ch05-有界胎紧紧}（另见\cref{remark:ch05-polish}），对任意序列 $\{\nu_k\}\subseteq N$，存在子列 $\{\nu_{k_j}\}\subseteq\{\nu_k\}$ 与 $\nu\in\mathscr{R}^+_f(X\times Y)$ 使
\begin{align}
\lim_{j\to\infty}\int_{X\times Y}g\,\mathrm{d}\nu_{k_j}
=\int_{X\times Y}g\,\mathrm{d}\nu
\quad\text{对所有 }g\in C_b(X\times Y).\label{align:ch05-5.65}
\end{align}
易见 $\nu\in\mathcal{Y}^+(X;Y)$。事实上，由于 $\{\nu_{k_j}\}\subseteq N\subseteq\mathcal{Y}^+(X;Y)$，由离散分解定理，等式\eqref{align:ch05-5.65}读作
\begin{align*}
\lim_{j\to\infty}\int_X\mathrm{d}\mu(x)\int_Y g(x,y)\,\mathrm{d}(\nu_{k_j})_x(y)
=\int_X\mathrm{d}\nu_{p_X}(x)\int_Y g(x,y)\,\mathrm{d}\nu_x(y)
\end{align*}
对所有 $g\in\mathcal{C}_b(X\times Y)$。在此等式中取 $g=g(x)$，得
\begin{align*}
\int_X g(x)\,\mathrm{d}\mu(x)=\int_X g(x)\,\mathrm{d}\nu_{p_X}(x)
\quad\text{对所有 }g\in C_b(X).
\end{align*}
由\cref{proposition:ch04-稠密}(ii)（另见\cref{proposition:appA-紧生成}），由此等式显然得 $\nu_{p_X}(E)=\mu(E)$ 对所有 $E\in\mathscr{B}(X)$。断言得证。

特别地，等式\eqref{align:ch05-5.65}对所有 $g\in\mathcal{C}_0(X\times Y)$ 成立。由于空间 $\mathcal{C}_0(X\times Y)$ 在 $L^1(X;C_0(Y))$ 中稠密（见\cref{proposition:ch04-稠密}(ii)与\cref{proposition:ch05-Caratheodory同构}），对任意 $h\in L^1(X;C_0(Y))$，存在序列 $\{g_m\}\subseteq\mathcal{C}_0(X\times Y)$ 使 $\lim_{m\to\infty}\|g_m-h\|_{L^1(X;C_0(Y))}=0$。固定 $\epsilon>0$，取 $m_0\in\mathbb{N}$ 充分大使 $\|g_{m_0}-h\|_{L^1(X;C_0(Y))}<\epsilon/3$。另一方面，由\eqref{align:ch05-5.65}存在 $j_0\in\mathbb{N}$ 使
\begin{align*}
\left|\int_{X\times Y}g_{m_0}\,\mathrm{d}\nu_{k_j}-\int_{X\times Y}g_{m_0}\,\mathrm{d}\nu\right|<\frac{\epsilon}{3}
\quad\text{对所有 }j>j_0.
\end{align*}
则对所有 $j>j_0$，
\begin{align*}
\left|\int_{X\times Y}h\,\mathrm{d}\nu_{k_j}-\int_{X\times Y}h\,\mathrm{d}\nu\right|
\leqslant\left|\int_{X\times Y}(h-g_{m_0})\,\mathrm{d}\nu_{k_j}\right|
+\left|\int_{X\times Y}(h-g_{m_0})\,\mathrm{d}\nu\right|
+\left|\int_{X\times Y}g_{m_0}\,\mathrm{d}\nu_{k_j}-\int_{X\times Y}g_{m_0}\,\mathrm{d}\nu\right|\\
<\left|\int_{X\times Y}(h-g_{m_0})\,\mathrm{d}\nu_{k_j}\right|
+\left|\int_{X\times Y}(h-g_{m_0})\,\mathrm{d}\nu\right|
+\frac{\epsilon}{3}.
\end{align*}
由于 $\{\nu_{k_j}\}\subseteq N\subseteq\mathcal{Y}^+(X;Y)$，由离散分解定理得
\begin{align*}
\left|\int_{X\times Y}(h-g_{m_0})\,\mathrm{d}\nu_{k_j}\right|
=\left|\int_X\mathrm{d}\mu(x)\int_Y(h-g_{m_0})(x,y)\,\mathrm{d}(\nu_{k_j})_x(y)\right|
\leqslant\|g_{m_0}-h\|_{L^1(X;C_0(Y))}<\frac{\epsilon}{3}.
\end{align*}
由于 $\nu\in\mathcal{Y}^+(X;Y)$，我们类似地得
\begin{align*}
\left|\int_{X\times Y}(h-g_{m_0})\,\mathrm{d}\nu\right|
=\left|\int_X\mathrm{d}\mu(x)\int_Y(h-g_{m_0})(x,y)\,\mathrm{d}\nu_x(y)\right|
\leqslant\|g_{m_0}-h\|_{L^1(X;C_0(Y))}<\frac{\epsilon}{3}.
\end{align*}

鉴于上述不等式，对任意 $\epsilon>0$，存在 $j_0\in\mathbb{N}$ 使对所有 $j>j_0$，
\begin{align*}
\left|\int_{X\times Y}h\,\mathrm{d}\nu_{k_j}-\int_{X\times Y}h\,\mathrm{d}\nu\right|<\epsilon
\quad\text{对所有 }h\in\mathcal{C}_0(X\times Y).
\end{align*}
综上，我们证明了：对任意序列 $\{\nu_k\}\subseteq N$，存在子列 $\{\nu_{k_j}\}\subseteq\{\nu_k\}$ 与 $\nu\in\mathcal{Y}^+(X;Y)$ 使 $\nu_{k_j}\stackrel{*}{\rightharpoonup}\nu$。由于 $N$ 胎紧，由\cref{proposition:ch05-Young胎紧一致}结论得证。
\end{proof}

如下结果类似地证明。

\begin{proposition}\label{proposition:ch05-逐点窄收敛}
设 $X$ 与 $Y$ 为具可数基的局部紧 Hausdorff 空间。设 $\{\nu_k\}\subseteq\mathcal{Y}^+(X;Y)$，$\nu\in\mathcal{Y}^+(X;Y)$ 满足对 $\mu$-几乎所有 $x\in X$，$(\nu_k)_x\stackrel{n}{\rightharpoonup}\nu_x$ 在 $\mathscr{R}^+_f(Y)$ 中。则 $\nu_k\stackrel{n}{\rightharpoonup}\nu$ 在 $\mathcal{Y}^+(X;Y)$ 中。
\end{proposition}

\begin{proof}
设 $h\in\mathcal{C}_b(X\times Y)$。由\cref{definition:ch05-Caratheodory}(ii)，对 $\mu$-几乎所有 $x\in X$ 有 $h(x,\cdot)\in C_b(Y)$。则对 $\mu$-几乎所有 $x\in X$，
\begin{align*}
\lim_{k\to\infty}\int_Y h(x,y)\,\mathrm{d}(\nu_k)_x(y)
=\int_Y h(x,y)\,\mathrm{d}\nu_x(y)
\end{align*}
（见\cref{definition:ch05-窄收敛}）。另一方面，对所有 $k\in\mathbb{N}$，
\begin{align*}
\left|\int_Y h(x,y)\,\mathrm{d}(\nu_k)_x(y)\right|
\leqslant\|h(x,\cdot)\|_\infty(\nu_k)_x(Y)=\|h(x,\cdot)\|_\infty,
\end{align*}
且由\cref{definition:ch05-Caratheodory}(ii)，映射 $x\mapsto\|h(x,\cdot)\|_\infty$ 属于 $L^1(X)$。于是由控制收敛定理，
\begin{align*}
\lim_{k\to\infty}\int_X\mathrm{d}\mu(x)\int_Y h(x,y)\,\mathrm{d}(\nu_k)_x(y)
=\int_X\mathrm{d}\mu(x)\int_Y h(x,y)\,\mathrm{d}\nu_x(y),
\end{align*}
结论得证。
\end{proof}




\end{document}